\section{Attack Methodology}
\label{sec:method}

\begin{figure}[h!]
\centering
\fbox{%
  \begin{minipage}[c][1.7in][c]{0.95\linewidth}
    \centering
    \textbf{Placeholder: ADSD Method Figure}\\[0.5ex]
    End-to-end attack pipeline. The final version should illustrate:
    (i) continuous prefix relaxation, (ii) projection-gap control with
    $\rho(U)$, (iii) verifier-aligned soft-prefix optimization, (iv)
    calibration-batch universal-prefix extension over target rollouts, and
    (v) discrete prefix evaluation in the speculative-decoding harness.
  \end{minipage}
}
\caption{Overview of the ADSD pipeline. Continuous prompt optimization is
coupled with explicit projection-gap control and verifier-aligned discrete
selection, then extended to a calibration-batch universal-prefix attack over
multiple prompts and timestamps.}
\label{fig:adsd-pipeline}
\end{figure}

\subsection{The Continuous-to-Discrete Projection Obstacle}

The attacker controls a discrete prefix, but gradient-based search is most
effective in a continuous relaxation of prompt space. Let
$E \in \mathbb{R}^{|\mathcal{V}| \times d}$ denote the shared token embedding
matrix. A relaxed prefix is
$U = (u_1, \dots, u_m)$, where each $u_t \in \Delta^{|\mathcal{V}|}$ is a
simplex vector over the vocabulary and the corresponding relaxed embedding is
\[
e_t(U) = E^\top u_t.
\]
This relaxation is attractive because it exposes gradient information over the
entire vocabulary at every position. The difficulty is that deployment still
requires a discrete projected prefix.

That obstacle is structural rather than incidental.

\begin{proposition}[Prompt Waywardness]
\label{prop:waywardness}
For any finite token embedding space with at least two tokens and any
$\delta > 0$, there exists a smooth objective $C$ over the convex hull of token
embeddings and a relaxed optimizer $U^*$ such that $U^*$ lies within distance
$\delta$ of a discrete token cell, yet the nearest-neighbor projection
$\Pi(U^*)$ is arbitrarily suboptimal among all discrete prefixes.
\end{proposition}

Prompt waywardness shows that proximity in embedding space does not imply
attack quality after projection. A successful continuous attack therefore
cannot rely on surrogate optimization alone; it must explicitly control the
relaxed-to-discrete gap.

\subsection{Bounding the Relaxation Gap}

Let $p_i^U$ and $q_i^U$ denote the target and draft next-token distributions at
speculative step $i$ under relaxed prefix $U$. We consider the weighted
disagreement objective
\[
C(U) = \sum_{i=1}^{\gamma} w_i \tv(p_i^U, q_i^U),
\]
where $w_i = \gamma - i + 1$ linearly emphasizes earlier speculative steps.
The corresponding discrete objective is
\[
C_d(z) := C(E(z)).
\]
To recover a deployable attack, define the token-wise nearest-neighbor
projection $\Pi(U) := (\Pi_1(U), \dots, \Pi_m(U))$ by
\[
\Pi_t(U) = \arg\min_{v \in \mathcal{V}} \|e_t(U) - E_v\|_2,
\]
and define the projection distortion
\[
\rho(U) := \sum_{t=1}^{m} \|e_t(U) - E_{\Pi_t(U)}\|_2.
\]

\begin{assumption}[TV-Lipschitzness]
\label{assump:tv-lipschitz}
For each speculative step $i$, there exist constants $L_i^p, L_i^q$ such that
for any two relaxed prefixes $U, V$,
\[
\tv(p_i^U, p_i^V) \le L_i^p \|U - V\|_{2,1},
\qquad
\tv(q_i^U, q_i^V) \le L_i^q \|U - V\|_{2,1},
\]
where
\[
\|U - V\|_{2,1} := \sum_{t=1}^{m} \|e_t(U) - e_t(V)\|_2.
\]
Let
\[
L_C := \sum_{i=1}^{\gamma} w_i (L_i^p + L_i^q)
\]
denote the aggregate Lipschitz constant of the weighted objective.
\end{assumption}

Under this assumption, the discrete quality of a projected relaxed solution can
be bounded explicitly.

\begin{theorem}[Relaxation Gap for Early Acceptance Collapse]
\label{thm:relax-gap}
Let $U^*$ be an $\varepsilon$-optimal solution of the relaxed problem such that
\[
C(U^*) \ge \sup_U C(U) - \varepsilon.
\]
Let
\[
z^* \in \arg\max_{z \in \mathcal{V}^m} C_d(z)
\]
be the optimal discrete prefix. Then
\[
C_d(\Pi(U^*)) \ge C_d(z^*) - \varepsilon - L_C \rho(U^*).
\]
\end{theorem}

\begin{proof}
By Assumption~\ref{assump:tv-lipschitz}, the objective $C$ is $L_C$-Lipschitz
under $\|\cdot\|_{2,1}$. Let $\hat{z} = \Pi(U^*)$. Since the continuous
embeddings of $\hat{z}$ are within $\rho(U^*)$ of $U^*$, we have
\[
C_d(\hat{z}) = C(E(\hat{z})) \ge C(U^*) - L_C \rho(U^*).
\]
Because the optimal discrete prefix is feasible in the relaxed space,
\[
C(U^*) \ge C_d(z^*) - \varepsilon.
\]
Combining the two inequalities yields the claim.
\end{proof}

We can also state an exact-recovery condition. Define the discrete attack
margin
\[
\Gamma := C_d(z^*) - \max_{z \neq z^*} C_d(z).
\]

\begin{corollary}[Exact Recovery under Discrete Margin]
\label{cor:exact-recovery}
Assume the optimal discrete adversarial prefix $z^*$ is unique
($\Gamma > 0$). If
\[
\varepsilon + L_C \rho(U^*) < \Gamma,
\]
then
\[
\Pi(U^*) = z^*.
\]
\end{corollary}

The theorem and corollary isolate the key design requirement: relaxed search
must improve the surrogate objective while also staying close to the discrete
token manifold. This is why ADSD uses an explicit projection-distortion
penalty rather than treating projection as a post-processing step.

\subsection{The ADSD Objective Formulation}

Motivated by Theorem~\ref{thm:relax-gap}, we construct a verifier-aligned
objective that couples disagreement maximization with explicit control of the
projection gap. The single-context attack targets the first verifier step,
since early rejection is the most damaging to token-wise speculative
throughput. Let $p^U$ and $q^U$ denote the target and draft next-token
distributions after the relaxed prefix $U$. ADSD optimizes
\begin{align*}
J(U) ={}&
\underbrace{\sum_{y} q^U(y)\,\mathrm{softplus}\!\left(\log q^U(y)-\log p^U(y)\right)}_{\text{soft collapse}}
+ \beta \,\tv(p^U, q^U)
+ \eta \,\mathrm{KL}(q^U \Vert p^U) \\
&- \tau \,\frac{1}{m}\sum_{t=1}^{m} H(u_t)
- \lambda \,\rho(U).
\end{align*}

The soft-collapse term rewards draft overconfidence on tokens that the target
assigns much smaller probability. The TV and reverse-KL terms enlarge the
surrounding disagreement region, the entropy term sharpens the relaxed simplex
vectors, and the projection penalty is the theory-guided component that keeps
the optimizer near a realizable discrete attack.

In the implementation, the prefix is parameterized by per-position vocabulary
logits. A softmax converts those logits into simplex vectors, and Adam updates
the logits directly. Each optimization step uses a straight-through argmax
estimator: the forward pass uses the hard projected token at each position,
while the backward pass differentiates through the softmax relaxation. The
released code therefore optimizes over the full vocabulary while keeping the
forward path faithful to the discrete token manifold used at deployment time.

Although the theory is phrased in terms of explicit nearest-neighbor
projection, the implementation uses argmax projection for computational
efficiency. The theoretical link is preserved by retaining the explicit
$L_2$ projection-distortion penalty $\rho(U)$ in the objective. Model
selection is verifier-aligned: although the optimizer follows the smooth
objective $J(U)$, the retained attack is ranked by its exact projected
first-token acceptance after hard projection in the actual
speculative-decoding pipeline.

In the experiments reported here, the core hyperparameters are a five-token
prefix, Adam learning rate $0.3$, temperature $1.0$, entropy weight $0.01$, TV
weight $1.0$, reverse-KL weight $0.1$, and gradient clipping at $1.0$. The
projection weight $\lambda$ is calibrated automatically from the ratio between
the unpenalized objective and the initial projection distortion, then swept on
Palmetto through a multiplicative factor. Because the search landscape is
highly non-convex, the implementation warm-starts the initial logits from a
short manual adversarial seed and searches the full vocabulary around that
initialization.

\subsection{Calibration-Batch Universal Prefix}

To move beyond a single boundary context, ADSD also optimizes a shared
universal prefix over a small calibration batch of GSM8K training prompts. For
each calibration prompt $x^{(j)}$, the target model first produces a fixed
continuation $y^{*(j)}_{1:K}$ of length $K$. These target-only rollouts define
the temporal contexts on which the universal prefix is optimized.

Let $p_i^{U,j}$ and $q_i^{U,j}$ denote the target and draft next-token
distributions at timestamp $i$ on calibration prompt $j$. The universal-prefix
objective is
\begin{align*}
J_{\mathrm{UAT}}(U) ={}&
\frac{1}{N}\sum_{j=1}^{N}
\Bigg[
\lambda_c \sum_{i=1}^{s} \bar{w}_i\,
\mathrm{SC}\!\left(q_i^{U,j}, p_i^{U,j}\right)
+ \beta \sum_{i=1}^{K} w_i\,\tv\!\left(q_i^{U,j}, p_i^{U,j}\right)
\\
&\qquad\qquad
+ \eta \sum_{i=1}^{K} w_i\,\mathrm{KL}\!\left(q_i^{U,j} \Vert p_i^{U,j}\right)
\Bigg]
- \tau \,\frac{1}{m}\sum_{t=1}^{m} H(u_t)
- \lambda \,\rho(U),
\end{align*}
where $w_i \propto K-i+1$ linearly emphasizes earlier timestamps and
$\bar{w}_i$ is the corresponding normalized restriction to the first $s$
timestamps. The term $\mathrm{SC}$ is the same soft-collapse signal used in the
single-context attack. The design preserves the theoretical emphasis on early
verifier failure while extending disagreement pressure across a full target
trajectory.

The universal-prefix implementation also includes a discrete initialization
stage before continuous refinement. Starting from the warm-start seed, ADSD
computes gradient-guided token replacements at each prefix position, retains
the top vocabulary candidates, and runs a small beam search over the prefix.
Candidate prefixes are ranked by exact projected first-token acceptance, with
weighted trajectory TV used as the tie-breaker. The best discrete candidate
then initializes the same straight-through optimization described above.

This hybrid discrete-plus-continuous pipeline is important in practice. It
places the continuous optimizer in a stronger basin, reduces dependence on a
single handwritten seed, and materially improves the quality of the projected
universal prefix recovered after refinement.
